\section{Conclusion}\label{sec:conclusion}

The central idea of Isegoria is to use opinion as a filter and evidence as the verdict, so that
the decision rests on what an item \emph{does} rather than on what the network \emph{thinks}. That
idea survives the analysis. Each level, however, has a property that the design documents did not
anticipate:
\begin{itemize}
  \item bridging scores are relative to their batch, and at the default penalties they retain most
    of the camp-size effect that bridging is meant to remove;
  \item latent-class DIF can mistake error in the ability proxy for a latent class;
  \item the evaluator's skill score is not proper, and the weight cap never binds;
  \item the anti-collusion discount depends on a detector that has almost no data within an epoch.
\end{itemize}
Each finding comes with a candidate correction. None of the corrections is validated yet.
The next steps are the evaluation plan of \cref{sec:plan}, which will turn provisional thresholds
into characterized ones, and independent review of both the statistics and the cryptography. The
paper is versioned with the repository and will be revised as those results arrive.

\section*{Acknowledgements and disclosure}
The analysis, the experiments and the text were prepared with the assistance of an AI system
(Claude, by Anthropic), used as a research and writing assistant. The author is responsible for
the content. Every numerical result can be regenerated with the scripts listed in
\cref{app:reproducibility}.
